% Appendix: Why K=3 with Continuous Rewards Outperforms K>3 with Binary Feedback
% Provides the detailed argument for why portfolio size is secondary to
% reward granularity in multi-model LLM routing.

\subsection{Reward Granularity and the \texorpdfstring{$K{>}2$}{K>2} Collapse}
\label{sec:reward_discrimination}

A natural question is whether ParetoBandit's $K{=}3$ portfolio is
a limitation relative to systems that evaluate larger pools.
PILOT~\cite{panda2025pilot} routes across $K{=}11$ models on
Routerbench (${\sim}36{,}500$ samples),
and PROTEUS~\cite{bhatti2026proteus} trains on $K{=}11$--$14$
portfolios.
We argue that \emph{reward granularity}, not portfolio size,
is the binding constraint for multi-model routing,
and that $K{>}2$ pools under binary rewards collapse to an
effective $K{=}2$ strong-vs.-weak triage.

\subsubsection{The Collapse Mechanism}

Modern LLMs at adjacent price tiers produce acceptable responses on
the majority of prompts.
A binary reward signal (correct/incorrect, win/loss, thumbs-up/down)
therefore clusters frontier and mid-tier models into a single
``strong'' category whenever both exceed the acceptance threshold
on the same prompt.
The multi-model routing problem structurally reduces to
$K{=}2$: one ``strong'' tier vs.\ one ``weak'' tier.

In our own $K{=}3$ portfolio, a natural binarisation threshold
($\tau{=}0.85$) causes the two stronger models (Gemini-2.5-Pro and
Mistral-Large) to tie on $85.8\%$ of prompts, eliminating the
router's ability to distinguish quality within a budget.
Switching from binary to continuous multi-dimensional rewards
(DeepSeek-R1 composite, $r_t \in [0,1]$) reduced the tied-prompt
rate from $85.8\%$ to ${<}1\%$ and raised best-arm entropy from
${\sim}1.1$ to ${\sim}2.7$ bits, restoring a genuinely multi-armed
problem.

\subsubsection{Evidence from PILOT (\texorpdfstring{$K{=}11$}{K=11}, Binary Rewards)}

PILOT's qualitative routing analysis (Section~5.1
of~\cite{panda2025pilot}) provides direct evidence of this
collapse at scale.
Despite having $K{=}11$ arms, per-task traffic concentrates on a
single model:

\begin{center}
\small
\begin{tabular}{@{}lll@{}}
\toprule
\textbf{Task} & \textbf{Dominant model} & \textbf{Share} \\
\midrule
MMLU          & GPT-4        & 90\% \\
ARC Challenge & GPT-4        & 89.4\% \\
GSM8K         & Claude-v1    & 94\% \\
MBPP          & GPT-4 + Claude & $\sim$72\% + 28\% \\
\bottomrule
\end{tabular}
\end{center}

\noindent
At most two models carry meaningful traffic on any given task;
the remaining 9--10 arms are effectively dead.
The appearance of diverse $K{=}11$ routing is an artefact of
\emph{aggregating across heterogeneous tasks}: within each task
the exact-match reward (binary 0/1) collapses the pool to
one or two active arms---precisely the $K{=}2$ triage
our analysis predicts.

\subsubsection{Corroboration from the Literature}

This structural limitation is widely reflected in the LLM routing
literature.
RouteLLM~\cite{ong2024routellm} and BaRP~\cite{wang2025barp}
both frame routing as an explicit binary strong/weak decision
trained on pairwise preference data; RouteLLM acknowledges the
$K{=}2$ restriction as a limitation and identifies extension to
multiple models as future work.
LLMRouterBench~\cite{li2026llmrouterbench} identifies
oversimplified binary accuracy metrics and imbalanced model
pools---where a single dominant model masks inter-tier
distinctions---as key evaluation failures for multi-model routing.

Conversely, every production-oriented router that supports
$K{>}2$ portfolios and demonstrates genuine multi-model
discrimination adopts continuous quality signals:
MixLLM~\cite{wang2025mixllm} uses predicted accuracy scores
across $K{=}5$ models, and
PROTEUS~\cite{bhatti2026proteus} optimises continuous SLA targets
across $K{=}11$--$14$ models via Lagrangian RL.
ParetoBandit follows this principle: our three-judge composite
reward ($r_t \in [0,1]$) produces inter-model gaps of
$0.793$--$0.932$ across the $K{=}3$ portfolio, which the router
exploits to route ${\sim}45\%$ of traffic at moderate budgets
to neither the cheapest nor the most expensive model---a
genuinely three-way decision.

\subsubsection{Implication for Portfolio Size}

The takeaway is not that large portfolios are undesirable, but
that \textbf{reward granularity is a prerequisite for portfolio
diversity to matter}.
A $K{=}11$ pool evaluated with binary exact-match rewards yields
at most $K{=}2$ effective routing decisions per task domain.
A $K{=}3$ pool evaluated with continuous multi-dimensional rewards
yields genuine three-way routing that varies with prompt context
and budget pressure.
Extending ParetoBandit to larger $K$ is straightforward (the
LinUCB update is $O(Kd^2)$ per step; the onboarding experiment
already demonstrates $K{=}3 \to K{=}4$ transitions), and the
continuous reward signal ensures that additional models will be
discriminated rather than collapsed into existing tiers.

We emphasise that this observation is specific to LLM routing,
where inter-model quality gaps are narrow and most prompts are
``easy'' for multiple tiers.
In bandit domains with larger inter-arm gaps, binary rewards
may carry sufficient separation.
